\section{A smaller cost for count crossings}
\label{v5cc:section}

The whole-count comparison contains a discrete Taylor remainder of the
binomial CDF. Its present bound pays separately for positive and negative
crossings, although those events are disjoint. It also uses a third-moment
estimate for crossings that can occur only beyond a fixed positive distance.
Both losses can be removed without changing the law or the comparison target.

Here is the coupling, including its endpoint convention. Let $U_{0,i}$ and
$U_{1,i}$, $1\le i\le n$, be independent centered residuals, with
$U_{0,i}\le a_0$, $U_{1,i}\ge-a_1$ and
$c=1-a_0-a_1\in(1/2,1]$. Put
\[
 V_j=\sum_{i\le j}U_{1,i}+\sum_{i>j}U_{0,i},\qquad0\le j\le n,
\]
and extend $V_j$ constantly outside $[0,n]$. The sequence $j+V_j$ is
strictly increasing, with increments at least $c$. For a legal count $k$
and $|u|\le1/2$, define
\[
 K=\max\{j\in\mathbb Z:j+V_j\le k+u\},\quad R=K-k,
 \quad t=\E V_k^2.
\]
If an independent selector has law $\operatorname{Bin}(n,p)$, its closed
event occurs exactly when the selector is at most $K$. Thus the mixture CDF
is $\E F_B(k+R)$. This construction uses the actual conditional laws; it
does not replace the convolution by independent choices of moment bounds.

Write $v=p(1-p)$, $\pi_j=\binom njp^j(1-p)^{n-j}$, with zero extension,
$m_B=\max_j\pi_j$, and $\chi=(n+1)p-k$. Assume $\chi<2$. For integer $r$,
the binomial difference identity and double summation give
\begin{equation}\label{v5cc:discrete}
 F_B(k+r)-F_B(k)-r\pi_k\le
 \frac{m_B}{(n+1)v}
 \left\{(\chi-1)_+r_++\chi_+\binom{r_-}{2}
                                      +\binom{r_-}{3}\right\}.
\end{equation}
For completeness,
$\pi_j-\pi_{j-1}=\pi_j((n+1)p-j)/(p(n-j+1))$.
When its numerator is positive, the denominator is at least $(n+1)v$.
For $r>0$, the positive differences have weights $r-i+1$ at $j=k+i$;
only $i=1$ can contribute positively because $\chi<2$.
For $r=-h$, the weights are $h-1-i$ at $j=k-i$, $0\le i\le h-2$.
The bound $(\chi+i)_+\le\chi_++i$ gives the two binomial coefficients
in \eqref{v5cc:discrete}. Zero extension handles counts outside $[0,n]$.

\begin{lemma}[Disjoint costs and a fourth-moment bound]
\label{v5cc:crossings}
Put $d=2c-1/2>0$ and
\[
 M_2(\chi,c,u)=\max\left\{
 \frac{(\chi-1)_+}{(c-u)^2},\,
 \frac{\chi_+}{[c+\min(u,c/3)]^2}\right\}.
\]
Then, pointwise,
\begin{align}
 (\chi-1)_+R_++\chi_+\binom{R_-}{2}
       &\le M_2(\chi,c,u)V_k^2,\\
 \binom{R_-}{3}&\le (V_k)_+^3/d^3,
 \qquad \binom{R_-}{3}\le (V_k)_+^4/d^4.
 \label{v5cc:fourth}
\end{align}
If also $|U_{s,i}|\le A$ and $t\le T$, then
\begin{equation}\label{v5cc:expected}
 \E\{F_B(k+R)-F_B(k)-R\pi_k\}
 \le\frac{m_Bt}{(n+1)v}
 \left[M_2(\chi,c,u)+
 \min\left\{\frac{\sqrt{A^2+3T}}{d^3},
                   \frac{A^2+3T}{d^4}\right\}\right].
\end{equation}
\end{lemma}
\begin{proof}
On $R=r\ge1$, crossing monotonicity gives $-V_k\ge cr-u$.
The difference
$(cr-u)^2-r(c-u)^2=(r-1)(c^2r-u^2)$ is nonnegative.
On $R=-h\le-2$, the count immediately above $K$ is outside the closed
event, so $V_k>c(h-1)+u$. Set $j=h-1\ge1$.
For $u\le c/3$, the function
$j(j+1)/(2(cj+u)^2)$ decreases for $j\ge1$: its derivative has the sign
of $j(2u-c)+u\le0$. Its maximum is $1/(c+u)^2$.
For $u>c/3$, lower $u$ to $c/3$ in this bound.
The positive and negative crossing events are disjoint. Hence the two
quadratic prices use a maximum, not their sum. When $R=0$ or $-1$,
both costs vanish.

For $h\ge3$, put $j=h-1\ge2$. Since $u\ge-1/2$,
\[
 \frac{\binom h3}{[c(h-1)+u]^3}
 \le\frac{j^3-j}{6(cj-1/2)^3}\le\frac1{(2c-1/2)^3}.
\]
For the last inequality the derivative numerator has the sign of
$2cj-(3j^2-1)/2<0$, using $c\le1$ and $j\ge2$; equality of the
last bound occurs at $j=2$. This proves the cubic estimate.
On this same event $V_k>d$, so $(V_k)_+^3/d^3\le(V_k)_+^4/d^4$.
For the other values of $R$, the left side is zero. Thus the quartic
estimate is pointwise too; no approximation at a crossing atom is used.

Independence and centering give, with variances $s_i$ of the $n$ terms,
\[
 \E V_k^4=\sum_i\E Z_i^4+6\sum_{i<j}s_is_j
 \le A^2t+3t^2\le t(A^2+3T).
\]
Cauchy--Schwarz also gives $\E(V_k)_+^3\le t\sqrt{A^2+3T}$.
Take expectations in the three pointwise estimates and then in
\eqref{v5cc:discrete}. The result includes $t=0$.
\end{proof}

The fourth-moment price is smaller whenever
$A^2+3T<(2c-1/2)^2$. It uses a moment bound already present in the
old proof. In that domain the square root is unnecessary for this term.
For larger noise retain the cubic price; taking the minimum preserves
both regimes. Dropping the individual fourth moments $A^2t$ would be
false, even for a centered three-point variable with rare nonzero values.

\subsection{The phase maximum still has three candidates}

In a closed scalar cell, $\chi\le b+u$. Define
\[
 g(u)=\max\left\{
 \frac{(b+u-1)_+}{(c-u)^2},\,
 \frac{(b+u)_+}{[c+\min(u,c/3)]^2}\right\},
 \qquad u_*=c-2b.
\]
For $[l,h]\subseteq[-1/2,1/2]$, its exact maximum is
\begin{equation}\label{v5cc:phase}
 \max\left\{g(l),g(h),
 \left.\frac1{4(c-b)}\;\right|\;
 c>b,\ l\le u_*\le h,\ u_*\le c/3\right\}.
\end{equation}
An ineligible third value is omitted. To prove this, the first summand
inside the maximum is nondecreasing: on its positive region its derivative
has positive numerator $c+u+2b-2>c-u>0$.
Below $c/3$, the positive part of the second has derivative
$(c-2b-u)/(c+u)^3$; above $c/3$ it is nondecreasing.
Its only possible interior maximum is therefore $u_*$, with value
$1/[4(c-b)]$. At that point $b+u_*=(c+u_*)/2\le3/4$, so the first
term vanishes. The supremum of a maximum is the maximum of the two
suprema. This proves \eqref{v5cc:phase}, including singleton intervals,
kinks and $c=1$.

For every allowed parameter, $g$ is bounded by the old sum of these two
terms. Both terms are nonincreasing in $c$. Consequently a conservative
cache at a smaller acquired $c$ is still valid, and the new exact phase
maximum is no larger than either the previous exact sum maximum or its
64-cell outward bound.

\subsection{What this changes in the local comparison}

Use the same acquired amplitude $A$, upper own-count variance $T$ and
lower gap coefficient $c$ as the current cell. Replace its curvature
coefficient by \eqref{v5cc:expected}, maximizing the phase by
\eqref{v5cc:phase}. All other terms, including the marked-crossing linear
term $\pi_k\E R$, the Gaussian repair and the actual binomial mass bound,
are retained. The new coefficient is never larger than the old one.
When passing from $1/n$ to $1/N$, clamp the complete improved inverse-index
correction at zero; its old strict positivity need not survive an improvement.
The existing negative event bracket is still required before using a lower
binomial mass bound.

For a fixed noise band, $A,T,c$ are frozen valid bounds. Hence this
replacement changes only the constant part of the inverse-index correction;
the existing convex endpoint argument in the noise parameter is unchanged.
It does not prove the acquisition of those bounds, cover all count cells,
or replace the separate zero-noise binomial theorem. It supplies a smaller
curvature contribution under exactly the old moment and coupling hypotheses.
